\documentclass[11pt]{article}

% Use packages
\usepackage[margin=1in]{geometry}
\usepackage{amsmath, amsfonts, amssymb}

% Title
\title{Linear Algebra Mastery \\ 
\large Vectors
}
\author{Evans N. Karago}
\date{\today}

% Main document
\begin{document}
\maketitle

% Introduction
\section{Introduction}
A vector is an \textit{ordered} list of finite real number. The list must have both magnitude and direction. Here are some examples of vectors:
\begin{enumerate}
\item Location and displacement
\item Colours 
\item Images
\item Video
\item Audio
\item Quantity
\item Daily return of stock
\item Stock portfolio and investment
\item Time series and signals
\item Word count and histogram
\item Cash flow in and out an entity or organization
\end{enumerate}
There are different ways to write a vector. In this article,  a vector will be a lower case letters with an arrow at the top(e.g. $\vec v$),  a  scalar will be represented by a lower case letter without and arrow at the top(e.g. $v$) or in special cases with a Greek letter (e.g. $\beta$). 
\section{Vector Operations}

% Vector addition
\subsection{Vector Addition}
We add vectors component by component. e.g.  $\vec u + \vec v$

\[
\begin{pmatrix}
u_1\\
u_2\\
\vdots \\
u_n
\end{pmatrix}
+
\begin{pmatrix}
v_1\\
v_2\\
\vdots\\
v_n
\end{pmatrix}
=
\begin{pmatrix}
u_1 + v_1\\
u_2 + v_2\\
\vdots\\
u_n + v_n
\end{pmatrix}
\]
Quick check:

With the following vectors $\vec v = (1, 3)$ and $\vec u = (2, -1)$, compute $\vec v + \vec u$.

\vspace{0.15in}
\begin{center}
\textbf{Solution:}
\end{center}
\[
\vec v + \vec u = \begin{pmatrix}
1 \\ 3
\end{pmatrix}
+ 
\begin{pmatrix}
2 \\ -1
\end{pmatrix}
= 
\begin{pmatrix}
3 \\
2
\end{pmatrix}
\]
\subsubsection*{Properties}
\begin{enumerate}
\item Commutative: $\vec v + \vec u $ = $\vec u + \vec v$
\item Associative: $\vec a ( \vec v + \vec u ) = (\vec a + \vec v) + \vec u$
\item Adding a zero vector to any vector has no effect e.g $\vec v + 0 = \vec v$
\item Subtracting a vector by itself yields a zero vector e.g. $\vec v - \vec v = 0`$
\end{enumerate}

% Scalar multiplication
\subsection{Scalar Multiplication}
A scalar is a \textit{real valued number}, sometimes coefficient of a vector. It's possible to scale a vector in ways we want, but before we get to that, below is a generalized way to scale a vector e.g. $\beta \vec v$

\[
\text{Let } \beta = c_1, c_2, \cdots, c_n \in \mathbb{R} \quad \text{and } \vec v  \in \mathbb{R}^n 
\]

\[
\beta(\vec v) = \beta \begin{pmatrix}
v_1\\
v_2\\
\vdots \\
v_n
\end{pmatrix}
=
\begin{pmatrix}
\beta v_1\\
\beta v_2\\
\vdots\\
\beta v_n
\end{pmatrix}
\]
Scaling a vector can lead to either of the following scenarios:
\begin{enumerate}
\item Stretching vector, this is true if $\beta > 1$
\item Shrinking vector, if $0 < \beta < 1$
\item Flipped vector, if $\beta = -1$
\item Flipped and stretched vector, if $\beta < -1$
\item Collapsed to zero vector, if $\beta = 0$
\end{enumerate}
Quick check:

Suppose $\vec v = (3, 2, 1)$, compute $2\vec v$.

\vspace{0.15in}

\begin{center}
\textbf{Solution: }
\end{center}
\[
2\begin{pmatrix}
3\\2\\1
\end{pmatrix}
= \begin{pmatrix}
6\\4\\2
\end{pmatrix}
\]

\subsubsection*{Properties}
\begin{enumerate}
\item Commutative: $\beta \vec v = \vec v \beta$, for some scalars $\beta$ and any vector $\vec v$
\item Associative: $(\beta \alpha)\vec v = \beta (\alpha \vec v)$, for some scalars $\beta, \alpha$ and any vector $\vec v$
\item Distributive: $(\beta + \alpha) \vec v = \beta \vec v + \alpha \vec v$
\end{enumerate}

% Linear combination
\section{Linear Combination}
Linear combination scales vectors and add them e.g. $c_1\vec v_1 + c_2\vec v_2 + \cdots + c_n \vec v_n$ give that, $c_i \in \mathbb{R}$ and $\vec v_i \in \mathbb{R}^n$. \\[0.1in]
Quick check: 

Given $\vec v_1 = (1, 2)$  and $\vec v_2 = (3, -1)$, compute $2\vec v_1 + 3 \vec v_2$

\begin{center}
\textbf{Solution:}
\end{center}

\[
2\cdot
\begin{pmatrix}
1 \\ 2
\end{pmatrix}
+
3 \cdot
\begin{pmatrix}
3\\-1
\end{pmatrix}
=
\begin{pmatrix}
2 \\4
\end{pmatrix}
+
\begin{pmatrix}
9 \\-3
\end{pmatrix}
=
\begin{pmatrix}
11\\ 1
\end{pmatrix}
\]
\subsubsection*{Special Combinations}
\begin{enumerate}
\item $c_ 1 = c_2 = c_3 = \cdots = c_n = 1$, sum of all vectors
\item $c_ 1 = c_2 = c_3 = \cdots = c_n = \frac{1}{n}$, Average of vectors.
\end{enumerate}
% Span
\section{Span}
Span is a set of all linear combinations of a set of vectors. e.g. 
\[
span\{\vec v_1, \vec v_2\} = \{c_1 \vec v_1, c_2 \vec v_2\mid c_i \in \mathbb{R}\}
\]
A span of two \textit{parallel} vectors in $\mathbb{R}^2$ is just a line, no new plane is formed, but if the vectors are not parallel, they will span the entire $\mathbb{R}^2$ plane.\\[0.1in]
Quick check:

Given $\vec v_ 1 = (1, 0)$ and $\vec v_2 = (0, 1)$, what is $span\{\vec v_1, \vec v_2\}$.

\vspace{0.15in}
\begin{center}
\textbf{Solution:}
\end{center}
\[
span\{\vec v_1, \vec v_2 \} = span\left\{c_1\begin{pmatrix}
1 \\0
\end{pmatrix},
c_2 \begin{pmatrix}
0 \\ 1
\end{pmatrix} \right\}\text{Given } c_1, c_ 2\in \mathbb{R}
\]
There are two special linear combinations of vectors yielding the sum of a set of vectors and average of the same.
\begin{enumerate}
\item Sum of vectors, $\beta_1 \vec v_1 = \beta_2 \vec v_2  = \cdots =\beta_n \vec v_n = 1$, for some vectors $\vec v_i \in \mathbb{R}^n$ and any scalars $\beta_i \in \mathbb{R}$.
\item Average of vectors, $\beta_1 \vec v_1 = \beta_2 \vec v_2 = \cdots = \beta_n \vec v_n = \frac{1}{n}$, for $\vec v_i \in \mathbb{R}^n  \text{ and } \beta_i \in \mathbb{R}$
\end{enumerate}
% Linear independence
\section{Linear Independence}
A collection or a list of n-vectors is said to be linearly independent if $\beta_1 \vec v_1 + \beta_2 \vec v_2 + \cdots + \beta_n \vec v_n = 0$, for all $\beta_i  = 0$.\\[0.1in]
Quick check:

Given $\vec v_1  = (1,2)$ and $\vec v_2 = (3, 5)$, determine if they are linearly independent.

\vspace{0.15in}
\begin{center}
\textbf{Solution:}
\end{center}
Recall: $\beta_1 \vec v_1 +\beta_2 \vec v_2+  \cdots + \beta_n \vec v_n = 0$
\[
\beta_1 \begin{pmatrix}
1 \\ 2
\end{pmatrix} + \beta_2 \begin{pmatrix}
3 \\ 5
\end{pmatrix}
= 0
\]

\begin{align}
 \beta_1 + 3 \beta_2 & = 0\\
 2 \beta_1 + 5 \beta_2 &= 0
\end{align}

\begin{align*}
2 \beta_1 + 6 \beta_2 & = 0\\
 2 \beta_1 + 5 \beta_2 &= 0\\
 \implies \beta_2 = 0 \text{ and } \beta_1 = 0 
\end{align*}
% Linear dependence
\section{Linear dependence}
A collection of n-vectors is linearly dependent if $\beta_1 \vec v_1 + \beta_2 \vec v_2 + \cdots + \beta_n \vec v_n = 0$, for some scalars $\beta_i$ not all zeros.\\[0.1in]
A collection of n-vectors containing the zero vector is always linearly dependent.\\[0.1in]
\textbf{Proof}:\\[0.1in]
Let $
S = \{\vec v_1, \vec v_2, \cdots \vec v_n\} \text{ and } \vec v_k  \in S \text{ such that } v_k = 0
$\\
Linear combination is given as 
\[\beta_1 \vec v_1 + \cdots + \beta_k \vec v_k + \cdots + \beta_n \vec v_n = 0
\]
Dividing by $\beta_k$
\[\frac{\beta_1}{\beta_k}\vec v_1 + \cdots + \vec v_k + \cdots + \frac{\beta_n}{\beta_k}\vec v_n = 0
\]
Finding $\vec v_k$
\[
\vec v_k =                                      -\left(\frac{\beta_1}{\beta_k}\right)\vec v_ 1 +  -\left(\frac{\beta_2}{\beta_k}\right)\vec v_ 2 +\cdots + -\left(\frac{\beta_{k-1}}{\beta_k}\right)\vec v_ {k -1} +-\left(\frac{\beta_{k + 1}}{\beta_k}\right)\vec v_ {k + 1} +\cdots +
-\left(\frac{\beta_n}{\beta_k}\right)\vec v_ n 
\]
From the equation above, it's evident that not all $\beta_i$ are all equal to zero, hence we can construct a new zero vector, which is redundant.
% Basis
\section{Basis}
A basis is a collection of n-vectors that are linearly independent and spans the entire space. Every vector in the space is reachable. Examples of a standard basis:
\[e_1 = \begin{pmatrix}
1 \\ 0 
\end{pmatrix} \qquad 
e_2 = 
\begin{pmatrix}
0 \\ 1
\end{pmatrix}
\]
\textbf{Proof:}\\[0.15in]
Suppose we have a collection of $n + 1$ n-vectors, given as 
$
S =
\{ 
\vec v_1, \vec v_2, \cdots , \vec v_n, \vec u \}$. The set  $S$ is linearly dependent, meaning that not all $\beta_i \neq 0$. \\[0.1in]
If we have
\[
\beta_1 \vec v_1 + \beta_2 \vec v_2 + \cdots + \beta_n \vec v_n + \beta_{n+1} \vec v_{n+1} = 0
\]
Then, $\beta_{n +1} = 0$, a contradiction for linear dependence.
% Machine learning application
\section{Machine Learning Application}
\begin{enumerate}
\item In machine learning, vectors can be represented by NumPy arrays.
\item Stretching, collapsing, or flipping a vector is a common practice while dealing with images, training neural network models.
\item Linear regression models, return their predictions as a linear combination of model weights and the target vector.
\item Multicollinearity causes a model prediction to be unstable. A collection of vectors that are linearly dependent causes multicollinearity.
\item Span tells us the true dimensionality of a set of vectors. A target vector should lie within the span, otherwise no exact solution will exist.
\end{enumerate}

\end{document}